\chapter{Grafy na powierzchniach, grafy dualne i homologia komórkowa}
\label{ch:tg-powierzchnie}

\section{Zanurzenia komórkowe i rodzaj powierzchni}
\label{sec:tg-cellular-surface}

\begin{definicja}[Zanurzenie komórkowe]
Niech $\Sigma$ będzie zwartą, spójną powierzchnią bez brzegu,
a $G$ skończonym spójnym grafem. Zanurzenie
$\iota:|G|\hookrightarrow\Sigma$ nazywamy \emph{komórkowym},
jeżeli każda składowa dopełnienia $\Sigma\setminus\iota(|G|)$
jest homeomorficzna z otwartym dyskiem. Te składowe nazywamy
\emph{ścianami} zanurzenia; po dołączeniu ich jako dwukomórek
otrzymujemy strukturę CW na $\Sigma$.
\end{definicja}

\begin{uwaga}[Dlaczego wymagana jest komórkowość]
Graf jednowierzchołkowy zanurzony w torusie ma dopełnienie
homeomorficzne z torusem bez punktu, a nie z otwartym dyskiem.
Nie wolno więc liczyć tego dopełnienia jako pojedynczej
dwukomórki. Równość $n-m+f=\chi(\Sigma)$ wymaga
zanurzenia komórkowego. W szczególności nie stosujemy jej
automatycznie do dowolnego rysunku grafu na torusie.
\end{uwaga}

\begin{twierdzenie}[Euler--Poincaré dla powierzchni orientowalnej]
\label{thm:tg-surface-euler}
Niech $\Sigma_g$ będzie zwartą, spójną, orientowalną powierzchnią
bez brzegu rodzaju $g\ge0$. Niech spójny graf $G$ będzie
zanurzony komórkowo w $\Sigma_g$, a $n,m,f$ oznaczają
odpowiednio liczby wierzchołków, krawędzi i ścian tego zanurzenia.
Wtedy
\[
 n-m+f=2-2g.
\]
\end{twierdzenie}

\begin{proof}
Z założenia wierzchołki, otwarte krawędzie i ściany
tworzą rozkład CW powierzchni $\Sigma_g$ na $n$
komórek wymiaru zero, $m$ komórek wymiaru jeden i $f$
komórek wymiaru dwa. Z definicji charakterystyki Eulera
$\chi(\Sigma_g)=n-m+f$. Klasyfikacja zwartych, spójnych,
orientowalnych powierzchni mówi, że
$\Sigma_g\cong\#^g(S^1\times S^1)$ i
$\chi(\Sigma_g)=2-2g$. Łącząc te równości,
otrzymujemy tezę. W tym dowodzie przyjmujemy znany,
wcześniej udowodniony w topologii powierzchni wzór
na charakterystykę Eulera \(\Sigma_g\); nie jest on
ukrytą konsekwencją samej definicji grafu.
\end{proof}

\begin{stwierdzenie}[Oszacowanie rodzaju z liczby krawędzi]
Niech prosty spójny graf $G$ o $n\ge3$ wierzchołkach
i $m$ krawędziach ma komórkowe zanurzenie w orientowalnej
powierzchni $\Sigma_g$. Wtedy
\[
 m\le3n-6+6g.
\]
Jeżeli $G$ jest dwudzielny, to
$m\le2n-4+4g$.
\end{stwierdzenie}

\begin{proof}
Każdy obchód ścianowy ma długość co najmniej trzy,
a w przypadku grafu dwudzielnego co najmniej cztery.
Istotnie, obchód długości jeden wymagałby pętli.
Obchód długości dwa w grafie prostym musiałby przejść
tę samą krawędź $uv$ w obu kierunkach.
Reguła następstwa półkrawędzi (opisana w kolejnej
sekcji) wymagałaby, aby zarówno przy $u$, jak i przy
$v$ była tylko jedna półkrawędź: jej cykliczny
następnik jest nią samą tylko przy stopniu jeden.
Spójny graf miałby zatem dokładnie dwa wierzchołki,
sprzecznie z $n\ge3$. W grafie dwudzielnym każdy
zamknięty spacer ma parzystą długość, więc dolne
ograniczenie wzrasta do czterech.
Każda krawędź ma dwie strony, liczone z krotnością
także wtedy, gdy należą do jednej ściany.
Stąd $3f\le2m$ lub $4f\le2m$. Podstawiając
$f=2-2g-n+m$, otrzymujemy odpowiednio
$3(2-2g-n+m)\le2m$ oraz $4(2-2g-n+m)\le2m$,
czyli żądane nierówności.
\end{proof}

\begin{przyklad}[Dlaczego $K_5$ potrzebuje uchwytu]
Dla $K_5$ mamy $n=5$ i $m=10$. Gdyby istniało
komórkowe zanurzenie w $\Sigma_0=S^2$, musiałaby zachodzić
nierówność $10\le9$, co jest niemożliwe.
Oszacowanie dla ogólnego rodzaju daje
$10\le9+6g$, a więc $g\ge1$.
Dla $K_{3,3}$ nierówność dwudzielna daje
$9\le 8+4g$, również więc $g\ge1$.
Jawne komórkowe zanurzenia obu grafów w torusie
skonstruujemy w przykładzie \ref{ex:tg-torus-rotations}.
W konsekwencji minimalny rodzaj orientowalny obu
grafów wynosi jeden. Samo oszacowanie daje jedynie
dolne ograniczenie.
\end{przyklad}

\section{Cykliczne porządki incydencji i grafy wstęgowe}
\label{sec:tg-rotation}

Samo abstrakcyjne $G=(V,E)$ nie zapamiętuje, jak krawędzie
są ułożone wokół wierzchołków na zorientowanej powierzchni.
Dane te można zakodować skończoną rodziną permutacji.

\begin{definicja}[Półkrawędź i układ rotacji]
Niech $G$ będzie skończonym spójnym multigrafem. Każda
krawędź daje dwie \emph{półkrawędzie} (ang. \emph{darts});
ich zbiór oznaczamy przez $D$. Także pętla daje dwie
różne półkrawędzie incydentne z tym samym wierzchołkiem.
Przez $\alpha:D\to D$ oznaczamy inwolucję bez punktów stałych,
która zamienia półkrawędzie tej samej krawędzi.
\emph{Układem rotacji} nazywamy permutację $\sigma$ zbioru $D$,
której cykle odpowiadają cyklicznej kolejności półkrawędzi
wokół poszczególnych wierzchołków. Dla wierzchołka
izolowanego kolejność jest pusta; w dalszych rachunkach
zakładamy graf spójny z co najmniej jedną krawędzią.
\end{definicja}

\begin{twierdzenie}[Powierzchnia wyznaczona przez układ rotacji]
\label{thm:tg-rotation-genus}
Niech $G$ będzie skończonym spójnym multigrafem o $n$
wierzchołkach i $m\ge1$ krawędziach, a $\alpha$ inwolucją parującą półkrawędzie i $\sigma$
permutacją, której cyklami są dokładnie zbiory półkrawędzi
przy kolejnych wierzchołkach, z zadanym porządkiem cyklicznym. Istnieje zorientowana, zamknięta,
spójna powierzchnia $\Sigma_g$ oraz komórkowe zanurzenie
$G\hookrightarrow\Sigma_g$ realizujące te porządki.
Jeżeli $f$ oznacza liczbę cykli permutacji
$\varphi:=\sigma\circ\alpha$ na $D$, to
\[
 2-2g=n-m+f,
 \qquad g=\frac{2-n+m-f}{2}.
\]
\end{twierdzenie}

\begin{proof}[Konstrukcja przez otoczenie wstęgowe]
Każdy wierzchołek zastępujemy małym zorientowanym
dyskiem. Na jego brzegu wyznaczamy rozłączne odcinki
odpowiadające incydentnym półkrawędziom, zachowując
kolejność cykliczną zapisaną w $\sigma$.
Dla każdej krawędzi doklejamy prostokąt
$[0,1]\times[0,1]$ do dwóch odpowiednich odcinków
brzegowych dysków. Doklejenia dobieramy w sposób
zachowujący orientację powstającej powierzchni z brzegiem.
Dopuszczamy, że oba końce prostokąta spotykają ten sam dysk,
co modeluje pętlę.

Uzyskujemy spójną, orientowalną powierzchnię wstęgową $N$
(z brzegiem), która deformacyjnie retraktuje się na $|G|$.
Śledzenie jej brzegu polega na przejściu wzdłuż danej
półkrawędzi na drugi koniec tej samej krawędzi (działanie
$\alpha$), a następnie do kolejnej półkrawędzi przy
wierzchołku (działanie $\sigma$). A więc składowe brzegu
$N$ odpowiadają dokładnie cyklom $\varphi=\sigma\alpha$.
Jest ich $f$.

Do każdej składowej brzegu $N$ doklejamy dysk.
Otrzymujemy zwartą, spójną powierzchnię zorientowaną
bez brzegu, w której $G$ jest szkieletem jednowymiarowym,
a każda ściana składa się z dołączonego dysku wraz
z przyległą częścią otoczenia wstęgowego poza szkieletem.
Istotnie, po usunięciu odcinków środkowych pasów
i promieni w dyskach każdy pas rozpada się na dwie
półwstęgi, a dyski na sektory pomiędzy kolejnymi
promieniami. Każdy sektor bez promieni i środka
dysku ma współrzędne kątową i radialną; także dla
jednego promienia używamy kąta w przedziale długości
$2\pi$. Łącząc sektory z półwstęgami zgodnie
z kolejnymi odcinkami brzegu, otrzymujemy dla
każdego obchodu kołnierz $S^1\times[0,1)$.
Okrąg $S^1\times\{0\}$ jest składową brzegu $N$,
a pominięty koniec kołnierza odpowiada szkieletowi.
Powtarzające się wierzchołki nie sklejają tych
kołnierzy, ponieważ punkty szkieletu usunięto.
Dołączenie dysku wzdłuż okręgu kołnierza daje
otwarty dysk, czyli ścianę komórkową.
Ponieważ $N$ ma charakterystykę $n-m$, po dołączeniu
$f$ dysków mamy $\chi(\Sigma)=n-m+f$.
Klasyfikacja powierzchni daje $\chi(\Sigma)=2-2g$,
co kończy dowód.
\end{proof}

\begin{przyklad}[Dwa układy rotacji tego samego grafu theta]
Rozważmy multigraf z wierzchołkami $u,v$ połączonymi trzema
równoległymi krawędziami. Niech $a_i$ i $b_i$ będą ich
półkrawędziami przy odpowiednio $u$ i $v$.
Wtedy $\alpha=(a_1b_1)(a_2b_2)(a_3b_3)$.
W obu modelach przy $u$ ustalamy porządek
$(a_1a_2a_3)$.

Jeśli przy $v$ przyjmiemy porządek $(b_1b_3b_2)$,
permutacja $\varphi=\sigma\alpha$ ma trzy cykle
$(a_1b_3)$, $(a_2b_1)$, $(a_3b_2)$.
Zatem $f=3$ oraz $2-2g=2-3+3=2$; otrzymujemy sferę.
Jeśli natomiast przy $v$ przyjmiemy ten sam porządek
$(b_1b_2b_3)$, $\varphi$ ma jeden cykl długości sześć:
$(a_1b_2a_3b_1a_2b_3)$. Teraz $f=1$ i
$2-2g=2-3+1=0$; otrzymujemy torus.
To samo abstrakcyjne $G$ ma więc dwa komórkowe zanurzenia
o różnych rodzajach powierzchni.
\end{przyklad}

\begin{center}
\begin{tikzpicture}[scale=.95,line cap=round]
  \begin{scope}[shift={(-2.9,0)}]
    \node[tgBlueVertex] (u) at (-1.05,0) {$u$};
    \node[tgRedVertex] (v) at (1.05,0) {$v$};
    \draw[tgBlue,line width=1.8pt] (u) to[bend left=43] (v);
    \draw[tgTeal,line width=1.8pt] (u)--(v);
    \draw[tgCoral,line width=1.8pt] (u) to[bend right=43] (v);
    \node[text=tgNavy,font=\small] at (0,-1.38) {$f=3$, $g=0$};
    \node[text=tgPurple,font=\footnotesize] at (0,-1.72) {$(b_1b_3b_2)$};
  \end{scope}
  \begin{scope}[shift={(2.9,0)}]
    \node[tgBlueVertex] (u) at (-1.05,0) {$u$};
    \node[tgRedVertex] (v) at (1.05,0) {$v$};
    \draw[tgBlue,line width=1.8pt] (u) to[bend left=43] (v);
    \draw[tgTeal,line width=1.8pt] (u)--(v);
    \draw[tgCoral,line width=1.8pt] (u) to[bend right=43] (v);
    \node[text=tgNavy,font=\small] at (0,-1.38) {$f=1$, $g=1$};
    \node[text=tgPurple,font=\footnotesize] at (0,-1.72) {$(b_1b_2b_3)$};
  \end{scope}
  \node[text=tgNavy,font=\footnotesize,align=center] at (0,-2.48)
    {Różnica leży w cyklicznym porządku \emph{półkrawędzi};\\
     dwuwymiarowy rysunek nie zastępuje danych rotacji.};
\end{tikzpicture}
\end{center}

\begin{przyklad}[Jawne zanurzenia $K_5$ i $K_{3,3}$ w torusie]
\label{ex:tg-torus-rotations}
W grafie prostym półkrawędź wychodzącą z $i$ ku $j$
oznaczamy parą $(i,j)$. Inwolucja $\alpha$ zamienia
$(i,j)$ na $(j,i)$, a $\sigma$ przesuwa drugi indeks
według porządku cyklicznego przy pierwszym indeksie.
W $K_5$ wybieramy następujące porządki sąsiadów:
\[
\begin{array}{c|ccccc}
i&1&2&3&4&5\\ \hline
\text{porządek}&(2345)&(1345)&(1254)&(1325)&(1423)
\end{array}
\]
Cykle $\sigma\alpha$ zapisujemy jako zamknięte obchody
wierzchołków (ostatni krok wraca do początku):
\[
\begin{gathered}
(1,2,3,5),\quad (2,1,3,2,4,5),\\
(3,1,4),\quad (4,1,5),\quad (4,2,5,3).
\end{gathered}
\]
Każdy krok $i\to j$ przechodzi w $j\to k$,
gdzie $k$ jest następnikiem $i$ w porządku przy $j$.
Obchody zawierają łącznie $4+6+3+3+4=20$ różnych
skierowanych półkrawędzi, czyli wszystkie.
Mamy $n=5,m=10,f=5$, stąd $2-2g=0$.

Dla $K_{3,3}$ z częściami $\{1,2,3\}$ i $\{4,5,6\}$
przy każdym wierzchołku pierwszej części wybieramy
porządek $(456)$, a przy każdym wierzchołku drugiej
porządek $(123)$. Obchody są następujące:
\[
\begin{gathered}
(1,4,2,5,3,6),\quad (4,1,5,2,6,3),\\
(5,1,6,2,4,3).
\end{gathered}
\]
Zawierają wszystkie $18$ półkrawędzi, więc
$n=6,m=9,f=3$ i ponownie $2-2g=0$.
Twierdzenie \ref{thm:tg-rotation-genus} daje w obu
przypadkach komórkowe zanurzenie w torusie.
\end{przyklad}

\section{Graf dualny na powierzchni}
\label{sec:tg-dual}

\begin{definicja}[Graf dualny]
Niech $G\subset\Sigma$ będzie skończonym spójnym grafem
zanurzonym komórkowo w zamkniętej, zorientowanej
powierzchni $\Sigma$, o zbiorze ścian $F$. W każdej ścianie wybieramy
punkt $v_F$. Dla każdej krawędzi $e$ grafu $G$ rysujemy
łuk $e^*$ łączący punkty dwóch ścian leżących po
przeciwnych stronach $e$, przecinający $e$ dokładnie
raz i nieprzecinający innych krawędzi pierwotnych.
Graf z wierzchołkami $v_F$ i krawędziami $e^*$
nazywamy \emph{grafem dualnym} i oznaczamy $G^*$.
Jeśli obie strony $e$ należą do tej samej ściany,
$e^*$ jest pętlą; dualny graf może mieć krawędzie równoległe.
\end{definicja}

\begin{twierdzenie}[Dualność komórkowa]
\label{thm:tg-cellular-dual}
Graf dualny można skonstruować jako graf zanurzony
komórkowo w tej samej powierzchni. Jest on spójny;
jego ściany odpowiadają wierzchołkom $G$.
Dla $m\ge1$ w oznaczeniach układu rotacji dual ma
parę permutacji $(\alpha,\sigma^*)$, gdzie
$\sigma^*=\sigma\alpha$ (przy odpowiednim oznaczeniu
półkrawędzi dualnych). Dwukrotne przejście do dualu
odtwarza układ rotacji grafu $G$.
\end{twierdzenie}

\begin{proof}
Przyjmijmy najpierw $m\ge1$. W modelu wstęgowym
do każdej ściany należy jeden dysk przyklejony
do okręgu brzegu otoczenia $N$. Dla każdego boku
pasa krawędzi wybieramy punkt na tym okręgu,
na wysokości środka pasa. W każdym dołączonym
dysku łączymy jego środek z wybranymi punktami
rozłącznymi (poza środkiem) promieniami.
W pasie łączymy dwa przeciwległe wybrane punkty
odcinkiem poprzecznym, który przecina krawędź
pierwotną dokładnie raz. W ten sposób otrzymujemy
wszystkie krawędzie dualne, bez przecięć między
ich wnętrzami. Jeżeli oba boki pasa należą do
jednej ściany, używamy dwóch różnych promieni
tego samego dysku; powstaje pętla dualna.

Rozcięcie pasów przez odcinki poprzeczne i dysków
ścianowych przez promienie pozostawia wokół każdego
wierzchołka pierwotnego obszar złożony z jego dysku
wierzchołkowego, sąsiadujących połówek pasów i sektorów
dysków ścianowych. Odcinając te części nieco przed
szkieletem dualnym, otrzymujemy dysk: do dysku
wierzchołkowego doklejamy kolejne fragmenty wzdłuż
przedziałów jego brzegu. Nieutożsamione brzegi tych
fragmentów zbliżają się do szkieletu dualnego;
sam szkielet jest wyłączony z obszaru. Powiększając
tak otrzymane dyski aż do szkieletu dualnego,
uzyskujemy otwarte dyski, po jednym przy każdym
wierzchołku pierwotnym. Są to wszystkie ściany dualne.

Równoważnie opiszmy kolejności cykliczne. Półkrawędź
dualną oznaczamy półkrawędzią pierwotną odpowiadającą
bokowi pasa, z którego wychodzi promień dualny.
Dwa boki jednego pasa paruje $\alpha$.
Przechodząc po okręgu ściany, mijamy kolejne boki
w porządku $\varphi=\sigma\alpha$, więc przy
zgodnym wyborze konwencji orientacji dostajemy
$\sigma^*=\varphi$. Obchody dualne są cyklami
$\sigma^*\alpha=\sigma\alpha^2=\sigma$, czyli
odpowiadają pierwotnym wierzchołkom. To samo
obliczenie daje $\sigma^{**}=\sigma$.
Spójność dualu wynika również z tego, że orbity
działań $\alpha,\sigma\alpha$ są dokładnie orbitami
działań $\alpha,\sigma$ (każda z tych par generuje
te same permutacje). Spójność $G$ oznacza jedną orbitę
na zbiorze półkrawędzi, więc dual jest spójny.

Jeżeli $m=0$, spójny graf ma jeden wierzchołek.
Komórkowość oznacza, że dopełnienie jest otwartym
dyskiem; powierzchnia jest sferą (jednopunktowe
uzwarcenie otwartego dysku). Dual ma jeden wierzchołek,
nie ma krawędzi i również jest komórkowo zanurzony.
\end{proof}

\begin{stwierdzenie}[Podstawowe własności dualności]
\label{thm:tg-dual-basic}
Dla skończonego spójnego grafu płaskiego $G\subset S^2$
graf dualny $G^*$ jest spójny i ma
$|V(G^*)|=f(G)$ oraz $|E(G^*)|=|E(G)|$.
Ponadto $(G^*)^*$ jest równoważny $G$ jako graf zanurzony
w sferze. Krawędź $e\in E(G)$ jest mostem wtedy i tylko
wtedy, gdy $e^*$ jest pętlą w $G^*$.
\end{stwierdzenie}

\begin{proof}
Pierwsze dwie równości wynikają z konstrukcji.
Aby wykazać spójność, bierzemy dwa punkty w różnych
ścianach i łączymy je krzywą na sferze, omijającą
wierzchołki pierwotnego grafu i przecinającą jego
krawędzie poprzecznie. Zapisując kolejne napotkane
ściany, otrzymujemy spacer w $G^*$.

Komórkowość i odtworzenie układu rotacji po
dwukrotnej dualizacji wynikają z twierdzenia
\ref{thm:tg-cellular-dual}. Na sferze układ rotacji
wyznacza osadzenie z dokładnością do zachowującego
orientację homeomorfizmu: odpowiadające sobie dyski
wierzchołkowe i pasy utożsamiamy zgodnie z rotacjami,
a homeomorfizmy okręgów brzegowych przedłużamy
na doklejone dyski. Daje to równoważność zanurzeń.

Jeżeli $e$ jest mostem, każdy zamknięty spacer przechodzący przez $e$ musi przejść przez nią także w przeciwnym kierunku: po usunięciu $e$ jej końce leżą w różnych składowych. Obchód ściany zawierający jeden bok $e$ zawiera zatem również drugi bok (każda półkrawędź występuje dokładnie raz wśród wszystkich obchodów). Stąd $e^*$ jest pętlą.
Odwrotnie, jeśli oba boki $e$ są po tej samej stronie,
to $e$ nie może należeć do żadnego cyklu prostego:
cykl tworzyłby krzywą Jordana, której dwa boki
przy $e$ należą do różnych ścian. Z kryterium mostu
wynika wtedy, że $e$ jest mostem.
\end{proof}

\begin{center}
\begin{tikzpicture}[scale=.88,line cap=round,line join=round]
 \begin{scope}[shift={(-3.1,0)}]
  \coordinate (a) at (-1.3,-1);\coordinate (b) at (1.3,-1);
  \coordinate (c) at (1.3,1);\coordinate (d) at (-1.3,1);
  \fill[tgBlue!9] (a)--(b)--(c)--cycle;
  \fill[tgTeal!9] (a)--(c)--(d)--cycle;
  \draw[tgEdge] (a)--(b)--(c)--(d)--cycle;
  \draw[tgAccent,line width=2pt] (a)--(c);
  \foreach \pt in {a,b,c,d} \node[tgBlueVertex] at (\pt) {};
  \node[text=tgBlue] at (.52,-.44) {$F_1$};
  \node[text=tgTeal] at (-.5,.45) {$F_2$};
  \node[text=tgNavy,font=\footnotesize] at (0,-1.55) {$G$};
 \end{scope}
 \begin{scope}[shift={(3.1,0)}]
  \node[tgBlueVertex] (f1) at (0,-.6) {$F_1$};
  \node[tgTealVertex] (f2) at (0,.8) {$F_2$};
  \node[tgRedVertex] (f3) at (1.8,.1) {$F_3$};
  \draw[tgAccent] (f1)--(f2);
  \draw[tgBlue,line width=1.7pt] (f1) to[bend right=27] (f3);
  \draw[tgBlue,line width=1.7pt] (f1) to[bend left=27] (f3);
  \draw[tgTeal,line width=1.7pt] (f2) to[bend left=27] (f3);
  \draw[tgTeal,line width=1.7pt] (f2) to[bend right=27] (f3);
  \node[text=tgNavy,font=\footnotesize] at (.7,-1.55) {$G^*$};
 \end{scope}
\end{tikzpicture}

{\footnotesize Dualność zamienia ściany w wierzchołki i zachowuje
liczbę krawędzi. Dwie krawędzie między tymi samymi ścianami
stają się równoległymi krawędziami dualnymi.}
\end{center}

\begin{twierdzenie}[Drzewo dualne dopełniające]
\label{thm:tg-dual-spanning}
Niech $G\subset S^2$ będzie spójnym grafem płaskim,
a $T$ jego drzewem rozpinającym. Wtedy zbiór
$\{e^*:e\in E(G)\setminus E(T)\}$ tworzy drzewo
rozpinające w $G^*$.
\end{twierdzenie}

\begin{proof}
Z twierdzenia Eulera $f=m-n+2$, a więc liczba krawędzi
dualnych, które pozostawiamy, wynosi
$m-(n-1)=f-1=|V(G^*)|-1$.
Wystarczy zatem wykazać spójność wybranego podgrafu dualnego.
Gdyby nie był spójny, istniałby niepusty przekrój
$\delta_{G^*}(S)$, którego żadna krawędź nie została
wybrana. Oznaczałoby to, że wszystkie jego krawędzie
pierwotne należą do $T$. Jednak przekrój w grafie
dualnym odpowiada zbiorowi cykli w grafie pierwotnym
(najprostszy przypadek to obwód brzegu sumy ścian
należących do $S$). Niech $A$ będzie zbiorem pierwotnych krawędzi odpowiadających
temu przekrojowi. Każdy wierzchołek pierwotny ma parzystą
liczbę incydencji z $A$: kiedy przechodzimy wokoło niego
po kolejnych sektorach, liczba zmian między ścianami z $S$
i spoza $S$ jest parzysta. Dotyczy to również pętli,
których dwie półkrawędzie liczymy osobno. Zatem $A$
wyznacza cykl nad $\mathbb F_2$. Ponieważ przekrój jest
niepusty, $A\ne\varnothing$ i graf $(V,A)$ zawiera cykl
(w razie pętli jest to cykl długości jeden). Wszystkie
krawędzie $A$ należałyby do drzewa $T$, sprzeczność.
Wybrany podgraf dualny jest więc spójny i ma $f-1$
krawędzi, czyli jest drzewem rozpinającym.
\end{proof}

\begin{uwaga}
Na powierzchni dodatniego rodzaju most nadal daje
pętlę dualną (dowód przez zamknięty spacer nie zależy
od rodzaju), ale implikacja odwrotna zawodzi.
W jednościanowym modelu torusa obie krawędzie
pierwotne są pętlami i nie są mostami, a ich duale
są również pętlami. Nie każdy cykl na torusie
rozdziela powierzchnię.
\end{uwaga}

\begin{twierdzenie}[Rozkład na drzewo i drzewo dualne]
\label{thm:tg-tree-cotree}
Niech $G$ będzie skończonym spójnym grafem komórkowo
zanurzonym w $\Sigma_g$. Dla każdego drzewa
rozpinającego $T$ grafu $G$ istnieje drzewo
rozpinające $T^*$ w grafie dualnym, złożone
wyłącznie z krawędzi dualnych do krawędzi spoza $T$.
Oznaczmy przez $B$ zbiór krawędzi pierwotnych
odpowiadających krawędziom $T^*$, a przez
$L=E(G)\setminus(E(T)\cup B)$ pozostałe krawędzie.
Wówczas $|L|=2g$.
\end{twierdzenie}

\begin{proof}
Rozpatrzmy podgraf dualny zawierający wszystkie
wierzchołki dualne i wszystkie krawędzie dualne
do krawędzi spoza $T$. Gdyby był niespójny,
wybieramy zbiór $S$ wierzchołków jednej składowej.
W pełnym grafie dualnym istnieje niepusty przekrój
między $S$ i jego dopełnieniem, gdyż dual jest spójny.
Każda krawędź tego przekroju odpowiada krawędzi $T$.
Przy każdym pierwotnym wierzchołku liczba przejść
między sektorami ścian z $S$ i spoza $S$, przy
jednym pełnym obrocie, jest parzysta.
Krawędzie pierwotne odpowiadające przekrojowi
tworzą więc niepusty podgraf drzewa $T$, w którym
wszystkie stopnie są parzyste. Jest to niemożliwe:
każdy niepusty zbiór krawędzi lasu ma składową
z krawędzią, a ta ma wierzchołek stopnia jeden
(na przykład koniec najdłuższej ścieżki).
Podgraf dualny jest zatem spójny i ma drzewo
rozpinające $T^*$.
Z rozłączności zbiorów krawędzi mamy
$|L|=m-(n-1)-(f-1)=2g$ na mocy
twierdzenia \ref{thm:tg-surface-euler}.
Dla $g=0$ otrzymujemy dokładnie twierdzenie
\ref{thm:tg-dual-spanning}; dla $g>0$
pozostaje $2g$ krawędzi, które wykorzystamy
do konstrukcji bazy homologii.
\end{proof}

\section{Od ścian do kompleksu łańcuchowego powierzchni}
\label{sec:tg-surface-chains}

\begin{definicja}[Operator brzegu dwukomórek]
Niech skończony spójny graf $G$ będzie komórkowo zanurzony
w zorientowanej zamkniętej powierzchni $\Sigma$.
Oznaczmy przez $F$ zbiór ścian. Definiujemy grupy
$C_2:=\mathbb Z^F$, $C_1:=\mathbb Z^E$ i
$C_0:=\mathbb Z^V$. Operator $\partial_1:C_1\to C_0$
jest macierzą incydencji z rozdziału \ref{ch:tg-algebra},
z zerową kolumną dla każdej pętli i osobną kolumną
dla każdej krawędzi równoległej.
Ustalamy dowolną orientację każdej krawędzi, a ściany orientujemy zgodnie z orientacją powierzchni,
a $\partial_2:C_2\to C_1$ wysyła zorientowaną ścianę
na sumę zorientowanych krawędzi jej obchodu brzegowego,
ze znakami zależnymi od kierunku przejścia.
\end{definicja}

\begin{stwierdzenie}[Warunek łańcuchowy]
Dla operatorów $\partial_2$ i $\partial_1$ określonych
powyżej zachodzi $\partial_1\partial_2=0$.
\end{stwierdzenie}

\begin{proof}
Każda ściana ma obchód będący zamkniętym spacerem.
Pod zastosowaniem $\partial_1$ wkład przy każdym kroku
jest różnicą końca i początku krawędzi. W sumie
wierzchołek będący końcem jednego kroku jest zarazem
początkiem następnego; wkłady teleskopowo się znoszą.
Zatem brzeg brzegu każdej komórki wymiaru dwa jest zerowy,
a przez liniowość zachodzi $\partial_1\partial_2=0$.
\end{proof}

\begin{twierdzenie}[Homologia komórkowa powierzchni orientowalnej]
\label{thm:tg-surface-homology}
Niech $G$ będzie skończonym spójnym grafem z komórkowym
zanurzeniem w zorientowanej, spójnej zamkniętej powierzchni
$\Sigma_g$. Po zmianie współczynników z $\mathbb Z$ na $\mathbb R$,
czyli po zastosowaniu $C_p\otimes_{\mathbb Z}\mathbb R$
i $\partial_p\otimes\operatorname{id}_{\mathbb R}$,
otrzymujemy
\[
H_0(\Sigma_g;\mathbb R)\cong\mathbb R,
\qquad H_1(\Sigma_g;\mathbb R)\cong\mathbb R^{2g},
\qquad H_2(\Sigma_g;\mathbb R)\cong\mathbb R.
\]
\end{twierdzenie}

\begin{proof}
Graf jest spójny, zatem $\operatorname{rank}\partial_1=n-1$
i $\dim\ker\partial_1=m-n+1$. Udowodnimy, że
$\ker\partial_2$ jest jednowymiarowe. W orientowanej
powierzchni każda krawędź sąsiaduje z dwiema stronami
ścian, z przeciwnymi znakami w ich obwodach. Oznaczmy zbiór ścian przez $\mathcal F$. Jeśli
$z=\sum_{F\in\mathcal F}a_F[F]$ spełnia
$\partial_2z=0$, to po porównaniu współczynnika
przy dowolnej krawędzi między ścianami $F,F'$ dostajemy
$a_F=a_{F'}$. Dualny graf ścian jest spójny:
dowolne dwie ściany można połączyć łukiem przecinającym
graf transwersalnie i omijającym wierzchołki.
Zatem wszystkie $a_F$ są równe. Z drugiej strony
suma wszystkich zorientowanych ścian ma brzeg zero,
bo każda krawędź pojawia się dwukrotnie z przeciwnymi
znakami. Stąd
$\ker\partial_2=\mathbb R\sum_F[F]$, a więc
$\operatorname{rank}\partial_2=f-1$.

Wynika z tego
\[
\dim H_1=\dim\ker\partial_1-\operatorname{rank}\partial_2
=(m-n+1)-(f-1)=m-n-f+2=2g,
\]
gdzie ostatnia równość wynika z twierdzenia Eulera--Poincarégo.
Ponadto $H_2=\ker\partial_2\cong\mathbb R$,
a $H_0=C_0/\operatorname{im}\partial_1\cong\mathbb R$.
To wyprowadza homologiczne znaczenie rodzaju powierzchni
bez pomijania roli operatora $\partial_2$.
\end{proof}

\begin{twierdzenie}[Homologia całkowita i baza cykli]
\label{thm:tg-surface-integral}
Przy oznaczeniach twierdzenia \ref{thm:tg-tree-cotree}
zachodzą izomorfizmy
\[
H_0(\Sigma_g;\mathbb Z)\cong\mathbb Z,\quad
H_1(\Sigma_g;\mathbb Z)\cong\mathbb Z^{2g},\quad
H_2(\Sigma_g;\mathbb Z)\cong\mathbb Z.
\]
Dla $e\notin E(T)$ niech $z_e$ będzie całkowitym
cyklem fundamentalnym: suma zorientowanej krawędzi
$e$ i skierowanej ścieżki w $T$ od końca do początku
$e$, ze znakami zgodnymi z orientacją krawędzi.
W przypadku pętli ścieżka ma długość zero.
Klasy $[z_e]$ dla $e\in L$ tworzą bazę
$H_1(\Sigma_g;\mathbb Z)$.
\end{twierdzenie}

\begin{proof}
Najpierw rozpatrzmy grupę cykli $Z_1=\ker\partial_1$.
Odwzorowanie $P:Z_1\to\mathbb Z^{E\setminus E(T)}$
zachowuje współczynniki przy krawędziach spoza $T$.
Każdy $z_e$ ma współczynnik $1$ przy $e$ i zero
przy pozostałych krawędziach spoza $T$, więc $P$
jest surjekcją. Jeżeli $P(z)=0$, to $z$ jest
łańcuchem wspartym na drzewie. Współczynnik
przy krawędzi incydentnej z liściem musi znikać,
gdyż $\partial_1z=0$. Usuwając liście po kolei,
otrzymujemy $z=0$. Zatem $P$ jest izomorfizmem.

Dla każdej krawędzi $e$ oznaczmy przez $F_+(e)$
i $F_-(e)$ ściany, których obchody przechodzą
przez $e$ odpowiednio zgodnie i przeciwnie
do jej orientacji. Mogą być tą samą ścianą.
Jeżeli $a=\sum_F a_F[F]$, współczynnik
$\partial_2 a$ przy $e$ jest równy
$a_{F_+(e)}-a_{F_-(e)}$. Niech
$Q:\mathbb Z^F\to\mathbb Z^B$ zachowuje tylko
współczynniki $\partial_2a$ przy krawędziach $B$.
Jest to operator różnic wzdłuż drzewa dualnego.
Dla dowolnej rodziny całkowitych liczb $(b_e)_{e\in B}$
wybieramy korzeń drzewa dualnego i nadajemy mu
wartość zero. Idąc od korzenia wzdłuż jedynej
ścieżki do każdej ściany, ustalamy kolejne $a_F$
przez dodanie lub odjęcie $b_e$, tak by zachodziło
$a_{F_+(e)}-a_{F_-(e)}=b_e$.
Nie występują sprzeczne warunki, bo drzewo
nie ma cykli. Otrzymujemy $Q(a)=b$.
Jądrem $Q$ są dokładnie stałe rodziny
współczynników, ponieważ drzewo jest spójne.
Każda stała rodzina daje $\partial_2 a=0$,
gdyż dwie strony każdej krawędzi mają
przeciwne znaki.

Weźmy dowolny cykl $z$. Wybieramy $a$ tak,
aby współczynniki $\partial_2a$ przy $B$
były równe współczynnikom $z$ przy $B$.
Cykl $z-\partial_2a$ reprezentuje tę samą klasę
homologii, a jego współczynniki przy $B$ są zerowe.
W bazie cykli fundamentalnych (izomorfizm $P$)
jest więc całkowitą kombinacją $z_e$ dla $e\in L$.
Jeżeli taka kombinacja jest brzegiem $\partial_2a$,
jej współczynniki przy $B$ są zerowe, więc $Q(a)=0$.
Rodzina $a_F$ jest stała, skąd $\partial_2a=0$.
Z niezależności cykli fundamentalnych wynika,
że wszystkie współczynniki kombinacji są zerowe.
Klasy $[z_e]$ tworzą zatem bazę nad $\mathbb Z$;
w szczególności nie ma torsji.

Ponadto $\ker\partial_2$ składa się ze stałych
rodzin, jest więc generowane przez sumę
wszystkich zorientowanych ścian. Ponieważ nie
ma trójkomórek, daje to $H_2\cong\mathbb Z$.
Obraz $\partial_1$ składa się dokładnie
z całkowitych łańcuchów zerowymiarowych o sumie
współczynników zero: każdy brzeg ma tę własność,
a różnicę dwóch wierzchołków otrzymujemy jako
brzeg skierowanej ścieżki między nimi.
Iloraz $C_0/\operatorname{im}\partial_1$
jest więc izomorficzny z $\mathbb Z$ przez
sumowanie współczynników. Korzystamy tu z twierdzenia
o homologii komórkowej, które utożsamia homologię
tego kompleksu z homologią singularną powierzchni.
\end{proof}

\begin{przyklad}[Sfera i torus -- co zmienia doklejenie ścian]
Rozważmy trójkąt jako graf na sferze. Ma jedną niezależną
klasę w $H_1(|G|;\mathbb Z)\cong\mathbb Z$, reprezentowaną
przez obieg trójkąta. Po dołączeniu dwóch ścian
(wewnętrznej i zewnętrznej) otrzymujemy sferę.
Brzeg pierwszej ściany jest dodatnim obiegiem trójkąta,
a drugiej jego przeciwnym obiegiem, więc
$H_1(S^2;\mathbb Z)=0$.

Standardowy model torusa ma z kolei jedną zerokomórkę,
dwie jednokomórki $a,b$ i jedną dwukomórkę dołączoną
wzdłuż słowa $aba^{-1}b^{-1}$. Operator $\partial_1$ znika,
a $\partial_2$ również znika, ponieważ suma wykładników
każdego generatora w słowie przyklejającym jest zerowa.
Stąd $H_1(T^2;\mathbb Z)\cong\mathbb Z^2$, choć grupa
fundamentalna samego szkieletu jednowymiarowego jest
wolna: $\pi_1(S^1\vee S^1)\cong F_2$.
Dołączenie dwukomórki dokłada relację
$aba^{-1}b^{-1}=1$, która zmienia grupę fundamentalną
na $\pi_1(T^2)\cong\mathbb Z^2$.
\end{przyklad}

\begin{center}
\begin{tikzpicture}[scale=.90,>=Stealth,line cap=round]
  \fill[tgBlue!8] (0,0) rectangle (4,2.7);
  \draw[tgEdge] (0,0) rectangle (4,2.7);
  \draw[tgAccent,-{Stealth[length=2.8mm]},line width=2.2pt] (.25,0)--(3.65,0);
  \draw[tgAccent,-{Stealth[length=2.8mm]},line width=2.2pt] (.25,2.7)--(3.65,2.7);
  \draw[tgTeal,-{Stealth[length=2.8mm]},line width=2.2pt] (0,.25)--(0,2.35);
  \draw[tgTeal,-{Stealth[length=2.8mm]},line width=2.2pt] (4,.25)--(4,2.35);
  \node[text=tgCoral,font=\small] at (2,-.35) {$a$};
  \node[text=tgCoral,font=\small] at (2,3.08) {$a$};
  \node[text=tgTeal,font=\small] at (-.36,1.35) {$b$};
  \node[text=tgTeal,font=\small] at (4.36,1.35) {$b$};
  \node[text=tgNavy,align=center] at (2,1.35) {$aba^{-1}b^{-1}$\\[2pt] jedna ściana};
\end{tikzpicture}

{\footnotesize Kwadrat fundamentalny torusa: przeciwległe boki
utożsamiamy zgodnie ze strzałkami, otrzymując $n=1$, $m=2$, $f=1$.}
\end{center}


% ============================================================
% Iteracja 4, rozdział 12
% ============================================================
